% EconPaperKit — 经济学论文 LaTeX 模板
% 支持中英文，适配经济学论文常用格式

\documentclass[12pt,a4paper]{article}

% ===== 中文支持 (取消注释以启用中文) =====
% \usepackage[UTF8]{ctex}
% \usepackage[heading]{ctex}

% ===== 页面布局 =====
\usepackage[margin=1in]{geometry}
\usepackage{setspace}
\onehalfspacing  % 1.5 倍行距

% ===== 数学与计量经济学 =====
\usepackage{amsmath,amssymb,amsfonts}
\usepackage{bm}           % 粗体数学符号
\usepackage{bbm}          % 黑板粗体

% 常用计量经济学符号
\newcommand{\E}{\mathbb{E}}
\newcommand{\Var}{\operatorname{Var}}
\newcommand{\Cov}{\operatorname{Cov}}
\newcommand{\plim}{\operatorname{plim}}
\newcommand{\indep}{\perp\!\!\!\perp}

% ===== 表格 =====
\usepackage{booktabs}     % 三线表
\usepackage{threeparttable}
\usepackage{dcolumn}      % 小数点对齐
\newcolumntype{d}[1]{D{.}{.}{#1}}

% 三线表回归结果宏
\newcommand{\regtable}[1]{%
  \begin{table}[htbp]
    \centering
    \caption{回归结果}
    \begin{tabular}{l d{4} d{4} d{4} d{4}}
      \toprule
      & \multicolumn{1}{c}{(1)} & \multicolumn{1}{c}{(2)} 
      & \multicolumn{1}{c}{(3)} & \multicolumn{1}{c}{(4)} \\
      \midrule
      #1
      \bottomrule
    \end{tabular}
  \end{table}
}

% ===== 图片 =====
\usepackage{graphicx}
\usepackage{tikz}
\usepackage{pgfplots}
\pgfplotsset{compat=1.18}

% ===== 参考文献 =====
\usepackage[backend=biber,style=authoryear,maxcitenames=2]{biblatex}
\addbibresource{references.bib}

% ===== 超链接 =====
\usepackage{hyperref}
\hypersetup{
  colorlinks=true,
  linkcolor=blue,
  citecolor=blue,
  urlcolor=blue
}

% ===== 其他 =====
\usepackage{enumitem}
\usepackage{appendix}
\usepackage[nottoc]{tocbibind}

% ===== 标题信息 =====
\title{论文标题\\\large 副标题（如有）}
\author{作者姓名\thanks{联系方式: email@example.com}}
\date{\today}

\begin{document}

\maketitle

% ===== 摘要 =====
\begin{abstract}
本文研究了[研究问题]，使用[数据来源]和[研究方法]，
发现[主要结论]。本文的贡献在于[学术贡献]。
\vspace{0.5cm}

\noindent\textbf{关键词:} 关键词1, 关键词2, 关键词3

\noindent\textbf{JEL 分类:} C01, C13, O40
\end{abstract}

% ===== 正文 =====
\section{引言}

[研究背景与意义]

本文的结构安排如下：第二节回顾相关文献；第三节建立理论框架；
第四节描述数据和变量；第五节报告实证结果；
第六节进行稳健性检验；第七节总结全文。

\section{文献综述}

[文献回顾与述评]

\section{理论框架}

\subsection{基准模型}

考虑如下计量模型：
\begin{equation}\label{eq:baseline}
  y_{it} = \alpha + \beta X_{it} + \gamma Z_{it} + \mu_i + \lambda_t + \varepsilon_{it}
\end{equation}

其中 $y_{it}$ 为被解释变量，$X_{it}$ 为核心解释变量，$Z_{it}$ 为控制变量，
$\mu_i$ 为个体固定效应，$\lambda_t$ 为时间固定效应。

\subsection{识别策略}

为处理潜在的内生性问题，本文采用[工具变量/双重差分/断点回归]方法。

\section{数据与变量}

\subsection{数据来源}

本文数据来自[数据来源]，样本期间为[起始年份]至[终止年份]，
涵盖[样本量]个观测值。

\subsection{变量定义}

\subsubsection{被解释变量}

[变量名]：衡量[经济含义]，数据来源为[来源]。

\subsubsection{核心解释变量}

[变量名]：衡量[经济含义]，数据来源为[来源]。

\subsubsection{控制变量}

包括[变量列表]，控制[潜在混淆因素]。

% === 描述性统计表 ===
\begin{table}[htbp]
  \centering
  \caption{描述性统计}
  \label{tab:desc}
  \begin{tabular}{l c c c c c}
    \toprule
    变量 & 观测数 & 均值 & 标准差 & 最小值 & 最大值 \\
    \midrule
    $y$   & 1,000 & 10.50 & 3.20 & 2.10 & 25.80 \\
    $X_1$ & 1,000 &  5.30 & 1.80 & 1.00 & 10.00 \\
    $X_2$ & 1,000 &  0.45 & 0.50 & 0.00 &  1.00 \\
    $Z_1$ & 1,000 & 20.10 & 8.50 & 5.00 & 45.00 \\
    \bottomrule
  \end{tabular}
  \begin{tablenotes}
    \item 注：此表报告了主要变量的描述性统计量。
  \end{tablenotes}
\end{table}

\section{实证结果}

\subsection{基准回归}

表~\ref{tab:baseline} 报告了基准回归结果。
第(1)列仅包含核心解释变量，第(2)列加入控制变量，
第(3)列加入个体固定效应，第(4)列进一步加入时间固定效应。

\begin{table}[htbp]
  \centering
  \caption{基准回归结果}
  \label{tab:baseline}
  \begin{tabular}{l d{4} d{4} d{4} d{4}}
    \toprule
    被解释变量: $y$ & \multicolumn{1}{c}{(1)} & \multicolumn{1}{c}{(2)} 
    & \multicolumn{1}{c}{(3)} & \multicolumn{1}{c}{(4)} \\
    \midrule
    $X_1$ & 0.523^{***} & 0.485^{***} & 0.412^{***} & 0.398^{***} \\
          & (0.082)     & (0.078)     & (0.065)     & (0.061)     \\
    $X_2$ & -0.234^{**} & -0.198^{**} & -0.176^{**} & -0.165^{**} \\
          & (0.095)     & (0.089)     & (0.072)     & (0.068)     \\
    \midrule
    控制变量  & No  & Yes & Yes & Yes \\
    个体 FE   & No  & No  & Yes & Yes \\
    时间 FE   & No  & No  & No  & Yes \\
    \midrule
    观测数    & 1,000 & 1,000 & 1,000 & 1,000 \\
    $R^2$     & 0.12  & 0.28  & 0.45  & 0.52  \\
    \bottomrule
  \end{tabular}
  \begin{tablenotes}
    \item 注：括号内为异方差稳健标准误。\\
    $^{*}p<0.10,\;^{**}p<0.05,\;^{***}p<0.01$
  \end{tablenotes}
\end{table}

\subsection{内生性处理}

[工具变量/2SLS/DID 结果]

\section{稳健性检验}

\subsection{替换被解释变量}

\subsection{替换样本期间}

\subsection{安慰剂检验}

\subsection{排除异常值}

\section{结论}

本文研究了[研究问题]，主要发现如下：
\begin{enumerate}
  \item [发现一]
  \item [发现二]
  \item [发现三]
\end{enumerate}

本文的政策含义是[政策建议]。研究的局限性在于[局限性]，
未来研究可以[未来方向]。

% ===== 参考文献 =====
\printbibliography

% ===== 附录 =====
\appendix
\section{变量定义表}
\section{附加回归结果}

\end{document}
